% ============================================================
% CAPITOLO 1 -- INTRODUZIONE
% ============================================================
\section{Introduzione}

L'elaborato ha come obiettivo la realizzazione di un'attività di verifica
software su un fork di Apache Avro, applicata a due classi del progetto.

La verifica è stata condotta attraverso diverse metodologie di testing:
test manuali basati sulla tecnica Category-Partition, test generati
automaticamente tramite Randoop, test generati con l'ausilio di un LLM,
e test guidati dalla copertura tramite EvoSuite.

Sono stati isolati i test relativi alle due classi in esame, escludendo
dall'esecuzione la suite di test già presente nel progetto originale.

\subsection{Apache Avro}

Apache Avro è un framework di serializzazione dati, nato all'interno del
progetto Hadoop. Fornisce un formato binario compatto ed efficiente per
la codifica dei dati, un sistema di schemi (definiti in formato JSON) che
descrivono la struttura dei dati serializzati, e un meccanismo di
\emph{schema evolution}, ovvero la possibilità di far comunicare tra loro
versioni diverse dello stesso schema senza rompere la compatibilità.
Quest'ultimo aspetto è direttamente rilevante per il progetto: una delle
due classi analizzate (\texttt{Resolver}) implementa proprio la logica di
risoluzione tra schema \emph{writer} e schema \emph{reader}.

\subsection{Classi selezionate}

La selezione delle classi è avvenuta applicando l'algoritmo di
\emph{class ranking} descritto nella consegna, basato sull'iniziale del
cognome dello studente applicata alla classifica dei code smell rilevati
da SonarCloud sul progetto Apache Avro.

Una prima classe candidata, è stata scartata per banalità, in accordo con l'indicazione della consegna di evitare
classi il cui testing risulti banale o impraticabile, è stata selezionata la classe successiva.

Le classi definitivamente selezionate sono:
\begin{itemize}
    \item \texttt{org.apache.avro.Resolver} (473 LOC, 25 code smell):
    classe di utilità che confronta, ricorsivamente, l'albero dello
    schema usato da chi ha \emph{scritto} i dati (\emph{writer}) con
    l'albero dello schema atteso da chi li deve \emph{leggere}
    (\emph{reader}). Per ogni coppia di nodi corrispondenti, determina
    quale azione è necessaria per adattare i dati scritti al formato
    atteso in lettura (nessuna azione, promozione di tipo, adattamento
    di record/enum/union, oppure segnalazione di incompatibilità),
    producendo un albero di oggetti \texttt{Resolver.Action}. È il
    meccanismo alla base della \emph{schema evolution} di Avro, ovvero
    la capacità di far comunicare correttamente produttori e
    consumatori di dati che utilizzano versioni diverse dello stesso
    schema;
    \item \texttt{org.apache.avro.util.Utf8} (172 LOC, 8 code smell):
    classe mutabile che rappresenta una sequenza di testo codificata
    direttamente come array di byte UTF-8, anziché come \texttt{String}
    Java (rappresentata internamente in UTF-16). Implementa le
    interfacce \texttt{CharSequence}, \texttt{Comparable} ed
    \texttt{Externalizable}, fornendo metodi per impostare e leggere il
    contenuto, confrontare istanze e calcolarne l'hash. È utilizzata dal
    modello dati generico di Avro per rappresentare i campi di tipo
    stringa evitando le conversioni di codifica ripetute che
    sarebbero necessarie utilizzando \texttt{String}.
\end{itemize}